\section*{Facilities, Equipment and Other Resources}

\subsection*{Georgia Institute of Technology}

The project will use Georgia Tech research and web infrastructure for
development, integration, and operation. CryoStack is currently deployed on a
Georgia Tech Linux virtual machine behind Nginx and supports HTTPS,
WebSocket-based applications, local-password/GitHub/ORCID authentication,
role-based administration, persistent experiment records, and connection to
remote computing resources. Institutional identity, networking, security,
backup, and server-administration services will be used as approved by Georgia
Tech. Ordinary faculty/student workstations support development, documentation,
testing, and laptop-scale tutorial workflows.

Georgia Tech's Partnership for an Advanced Computing Environment (PACE)
provides Slurm-managed high-performance computing, parallel filesystems,
software modules, and user support. PACE has already been used for ICESEE MPI
development and scaling experiments and will serve as the reference backend
for adapter qualification and science demonstrations. Specific allocations,
queues, storage quotas, and staff commitments will be confirmed before
submission and described without assigning a monetary value as voluntary
committed cost sharing.

\subsection*{Software and Deployment Resources}

Existing project resources include the ICESEE ensemble data-assimilation
framework; the CryoStack gateway, Control Center, and modular CryoLauncher web
application; the LIVIST radar/borehole englacial-temperature explorer; the
manifest-driven Frozen Legacies historical-radar catalog and its first LYRA
adapter; desktop HPC connectors for Linux, macOS, and Windows; SSH/Slurm
execution code; modular AWS identity, IAM, networking, storage, registry, and
Batch components; Docker/OCI and Apptainer
recipes; Spack packages and environments; model integrations for ISSM,
Icepack, flowline, and Lorenz-96; and public documentation, tests, and issue
tracking. The project also builds on PISM, PETSc, MPI, Firedrake, Jupyter Book,
GitHub, and Zenodo, subject to their licenses and terms.

AWS services will be used only for bounded backend engineering, qualification,
and demonstrations covered by the project budget, and for the bounded
computing-access pilot if it is approved. CryoStack does not require AWS;
outside these activities, cloud runs are billed by AWS to the researcher's own
account. PACE remains the reference environment, while the AWS work
establishes a documented path for project-managed and bring-your-own-account
execution.

\subsection*{Collaborative Resources}

The project will seek letters of collaboration documenting bounded technical
review or user-evaluation roles from appropriate PISM users/developers,
historical-radar data stewards, Georgia Tech PACE/web-security personnel, and
research-software or science-gateway specialists. These collaborations will advise interface, security, usability,
and portability decisions; they will not be represented as committed resources
until the corresponding letters and institutional approvals are obtained.
